% Generated draft: requires mathematical and source review.
\documentclass[11pt]{article}
\usepackage[margin=25mm]{geometry}
\usepackage{fontspec}
\setmainfont{Latin Modern Roman}
\usepackage{amsmath,amssymb,mathtools,unicode-math}
\setmathfont{Latin Modern Math}
\usepackage{enumitem,needspace,xurl,hyperref}
\setlength{\parindent}{0pt}
\setlength{\parskip}{5pt}
\setlength{\emergencystretch}{3em}
\newcommand{\N}{\mathbb{N}}
\newcommand{\Z}{\mathbb{Z}}
\newcommand{\Q}{\mathbb{Q}}
\newcommand{\R}{\mathbb{R}}
\newcommand{\C}{\mathbb{C}}
\newcommand{\F}{\mathbb{F}}
\newcommand{\A}{\mathbb{A}}
\newcommand{\PP}{\mathbb{P}}
\newcommand{\E}{\mathbb{E}}
\newcommand{\eps}{\varepsilon}
\newcommand{\dd}{\,\mathrm{d}}
\begin{document}
\section*{UnsolvedMath ID: 30000867}
\textbf{Source problem code:} catalogue entry
\subsection{Definitions and mathematical statement}
Let $E$ and $F$ be number fields and let $M$ be a motive over $F$ with coefficients in $E$. Let $\mathbb A_F$ denote the adele ring of $F$. The source requests a connected reductive algebraic group $G/F$, an automorphic representation $\pi$ of $G(\mathbb A_F)$, and a representation $r$ of the $L$-group ${}^LG$, such that \[L(M,s)=L(s,\pi,r).\] In the usual complex $L$-group convention, ${}^LG$ is built from the complex dual reductive group and the relevant Galois or Weil action, and $r$ is a finite-dimensional representation subject to the appropriate algebraicity/continuity conditions. These conventions and a coefficient embedding $E\hookrightarrow\C$ are not specified in the extract. The first task is an automorphic realization for the given motive; the second is to choose $G$ as small as possible among such realizations. No ordering by dimension, rank, inclusion, isogeny, or another size criterion is supplied, so the minimization cannot be given a unique precise meaning. Likewise the source does not specify a category of motives, purity/effectivity assumptions, the allowed automorphic spectrum, or whether the equality compares completed $L$-functions, all finite Euler factors, or only factors outside a finite set of places. No cuspidality, temperedness, splitness of $G$, or faithfulness of $r$ is inserted.\par The supplied metadata are internally discordant. The original and clean/current statements explicitly give the realization and minimal-group tasks, and the corrected\_verified note calls the problem explicit. In contrast, the 2026-08-21 literature assessment says the extracted sentence only records nonuniqueness of a parametrization, omits its question, and yields no standalone claim, with an open classification for that reason. That description does not match the supplied displayed formulation. Both are preserved rather than using the assessment to erase the explicit question or asserting that the discrepancy has been resolved.
\subsection{Short English statement}
Realize a motive's $L$-function as an automorphic $L$-function and seek the smallest possible reductive group. The motive category, equality of local/global factors, and meaning of smallest are unstated; the literature summary also conflicts with the supplied explicit question.
\subsection{Sources}
\begin{itemize}\item Local UnsolvedMath registry, permanent ID 30000867, source problem catalogue entry, \emph{Automorphic Realization of Motivic L-Functions}, proposed year 2007; no individual proposer supplied.\item Supplied citation: \emph{Modulformen} (2008), Oberwolfach Reports, DOI 10.4171/owr/2007/51, pp.\ 3013--3072: \url{https://doi.org/10.4171/owr/2007/51}. Supplied literature label: original Oberwolfach report.\end{itemize}
\end{document}
